\chapter{Conclusion}
\label{ch:conclusion}

This dissertation developed a hybrid detector for sensor attacks under incomplete observations. The experiments showed that strong individual experts do not necessarily make a better combined detector: their detections may overlap while their false positives accumulate. The final design addresses this problem by combining complementary evidence and checking the alarms added by specialists.

General combines five internal residual classifiers using a target-group-excluding reference. Drift and Noise add temporal evidence with a declared three-hour decision delay. Routing, feedback and verification check the specialist alarms while General retains an immediate decision path. This design developed through graph and temporal experiments, inspection of expert errors and the rejection of alternatives that failed the evaluation criteria.

The two networks benefited from different changes. In Modena, guarded integration raised pooled F1 from 0.6863 to 0.8672 and reduced false positives by 84.2\%, with small average losses in drift and noise F1. In L-Town, adding received-pressure history to General improved every family mean, pooled F1 and mean clean-period FPR in the final matched comparison.

The final assessment retained ten model seeds and six evaluation sources per network. Modena achieved mean pooled F1 of 0.8669 and family macro F1 of 0.9045, with all five family means above 0.80. L-Town achieved 0.8950 and 0.9083, respectively. Four of its family means exceeded 0.89; replay remained at 0.7312. Clean-period false-positive rates were 0.0584\% for Modena and 0.0550\% for L-Town. Additional fits preserved the broad performance profiles.

The similar macro scores conceal different weaknesses. Modena's drift and noise performance varies by source. In L-Town, replay remains difficult despite improving in every paired history comparison; the later 0.75 objective was not met. Performance on the other families remains strong in the larger network.

The separate three-seed reference factorial in Section~\ref{sec:reference_factorial} isolates iterative reweighting and target-group exclusion within the retained representation. After downstream refitting and common-policy calibration, iterative reweighting improves pooled F1 in both networks and substantially improves Modena replay. Exclusion supplies no consistent pooled-F1 benefit, and L-Town replay depends on the interaction between the two operations. These findings support a narrower account of the reference's contribution: iterative reweighting helps under the tested recipe, while the exclusion constraint has no consistent performance advantage.

At the original operating points and against the pooled-F1-calibrated control with the same maximum three-hour allowance, the hybrid improves pooled F1 and temporal-family coverage in Modena. In L-Town, the simpler classifier has slightly higher pooled F1 and fewer clean false alarms, while the hybrid has higher replay, drift and noise F1 in every paired cell. Supplementary pooled-F1 calibration reverses both L-Town orderings at positive calibration caps, while the zero cap favors the simpler control in both networks. The fixed-rule diagnostics show that specialists improve pooled F1 and their intended families, while verification limits background errors; they also expose a replay cost from L-Town's Noise alarms. The added complexity has conditional benefits, determined by the attack families and false-positive allowance.

These results concern simulated, observed pressure endpoints under fixed attack and missingness conditions. Repeated fits share evaluation sources, and specialist decisions retain their stated delay. The supplementary analyses reuse those sources and do not replace the original selected system or provide fresh confirmation. Incident-level evaluation, wider observation conditions and computational measurements are still needed before operational use.
